\section*{Chapitre \ref{ch:b2:speciation} --- Spéciation et macroévolution}

\begin{solution}{exo:b2:speciation:1}
Une espèce est un groupe de populations naturelles interfécondes,
isolées reproductivement des autres groupes de ce type. Le concept échoue
pour les lignées asexuées (bactéries, rotifères bdelloïdes, pissenlits~:
concept morphologique ou phylogénétique, groupes de formes semblables ou
lignées diagnosables)~; pour les fossiles (concept morphologique, puisque
aucun croisement n'est possible)~; et pour les populations allopatriques
qui ne se rencontrent jamais (concept écologique ou phylogénétique,
puisque l'interfécondité ne peut pas être testée) --- et il est brouillé
par les chênes qui s'hybrident et par les espèces en anneau.
\end{solution}

\begin{solution}{exo:b2:speciation:2}
Grenouilles à des mois différents~: prézygotique (temporel). Mulet~:
postzygotique (stérilité des hybrides). Spermatozoïde incapable de se
fixer~: prézygotique (gamétique). Hybrides vigoureux aux \omterm{def:b2:angiosperm-reproduction:seed}{graines}
stériles~: postzygotique (stérilité des hybrides, ou dégénérescence des
hybrides si l'échec survient à la génération suivante). Grillons aux
chants différents~: prézygotique (comportemental).
\end{solution}

\begin{solution}{exo:b2:speciation:3}
Allopatrique~: une barrière divise une aire de répartition (les
crevettes-pistolets des deux côtés de l'isthme de Panama). Péripatrique~:
quelques fondateurs au-delà de l'aire (les pinsons de Darwin issus d'un
ancêtre continental). Sympatrique~: sans barrière (la mouche de la pomme
sur l'aubépine et le pommier~; un blé \omterm{def:b2:genome-diversification:polyploidy}{allopolyploïde}). Parapatrique~: le
long d'un gradient, avec une \omterm{prop:b2:speciation:reinforcement}{zone hybride} au milieu (des graminées sur et
hors des déblais miniers, fleurissant à des époques différentes).
\end{solution}

\begin{solution}{exo:b2:speciation:4}
L'espèce 1 l'emporte toujours~; l'espèce 2 l'emporte toujours~;
coexistence stable~; l'une ou l'autre l'emporte selon les effectifs
initiaux. La coexistence est stable lorsque chaque espèce peut envahir
l'autre à sa capacité de charge, c'est-à-dire lorsque chaque espèce se
limite elle-même davantage qu'elle ne limite l'autre
($K_1 > \alpha K_2$ et $K_2 > \beta K_1$).
\end{solution}

\begin{solution}{exo:b2:speciation:5}
$\alpha K_2 = 400 < K_1 = 1000$~: l'espèce 1 envahit~; $\beta K_1 = 600
< K_2 = 800$~: l'espèce 2 envahit. $N_1^{*} = (1000 - 400)/(1 - 0.3) =
857$, $N_2^{*} = (800 - 600)/0.7 = 286$~; à l'équilibre, l'espèce
1 est en dessous de son $K$ de $\alpha N_2^{*} = 143$ et l'espèce 2 de $\beta
N_1^{*} = 514$.
\end{solution}

\begin{solution}{exo:b2:speciation:6}
$\alpha K_2 = 1200 > K_1$~: l'espèce 1 ne peut pas envahir l'espèce 2 à
sa capacité de charge~; $\beta K_1 = 300 < K_2$~: l'espèce 2 peut envahir
l'espèce 1. L'espèce 2 l'emporte toujours, parce que ses individus nuisent
à l'espèce 1 plus que ceux de l'espèce 1 elle-même, alors qu'elle-même
n'est guère affectée. Avec $\alpha = \beta = 1.5$~: $\alpha K_2 = 1200 > 1000$
et $\beta K_1 = 1500 > 800$~: aucune ne peut envahir l'autre, chacune
exclut l'autre une fois installée, et la première arrivée l'emporte.
\end{solution}

\begin{solution}{exo:b2:speciation:7}
Le tétraploïde $4n$ produit des gamètes $2n$~; croisé avec les gamètes $n$
du \omterm{def:b2:meiosis-heredity:meiosis}{diploïde}, il donne des triploïdes $3n$, dont la \omterm{def:b2:meiosis-heredity:meiosis}{méiose} ne peut pas
apparier trois lots et produit des gamètes aneuploïdes, pour la plupart
non viables. Le
tétraploïde ne peut se reproduire qu'avec ses semblables (ou par
autofécondation)~: il est isolé en une seule étape. Il échoue en général
parce qu'il est seul --- ses seuls partenaires possibles sont des
\omterm{def:b2:meiosis-heredity:meiosis}{diploïdes} et chacun de ces croisements est perdu --- et se trouve
submergé~; il persiste lorsqu'il s'autoféconde, se multiplie par voie
végétative ou apparaît de façon répétée.
\end{solution}

\begin{solution}{exo:b2:speciation:8}
$0.01\times 10\times 10 = 1$ incompatibilité~; avec 30 chacune, $0.01\times
900 = 9$. L'isolement croît comme le carré du temps de divergence~: deux
populations qui ont divergé trois fois plus longtemps sont neuf fois plus
isolées --- la boule de neige.
\end{solution}

\begin{solution}{exo:b2:speciation:9}
$w = 2\sqrt{8/0.05} = \qty{25}{km}$~; $w = 0.3\sqrt{8/0.4} =
\qty{1.3}{km}$. Le renforcement est plus probable dans la seconde~: les
hybrides sont fortement désavantagés, si bien qu'une femelle qui refuse
l'autre forme y gagne beaucoup, et la zone est assez étroite pour que les
deux formes se rencontrent assez souvent pour que ce refus soit
sélectionné.
\end{solution}

\begin{solution}{exo:b2:speciation:10}
$\alpha K_2 = 450 < 500$, si bien que chacune envahit l'autre~: elles
coexistent, à $N^{*} = (500 - 450)/(1 - 0.81) = 263$ chacune --- de
justesse, au bord de l'exclusion, et une petite variation de $K$ ferait
basculer l'issue. Avec
$\alpha = \beta = 0.4$~: $N^{*} = (500 - 200)/(1 - 0.16) = 357$ chacune.
Le déplacement a fait pivoter chaque isocline en l'éloignant de l'autre
($K/\alpha$ s'est déplacé vers l'extérieur), si bien que le croisement
s'est éloigné des axes et que chaque espèce se trouve plus près de son
propre $K$~: une coexistence plus sûre, avec davantage d'individus de
chaque espèce.
\end{solution}

\begin{solution}{exo:b2:speciation:11}
$R = h^{2}S = 0.7\times(-1) = \qty{-0.7}{mm}$~: comme observé. Lors de la
sécheresse de 1977, en l'absence de géospize à gros bec, les mêmes
géospizes à bec moyen furent sélectionnés pour des becs \emph{plus gros},
puisque les grosses \omterm{def:b2:angiosperm-reproduction:seed}{graines} étaient la nourriture restante~; en 2004, les
grosses \omterm{def:b2:angiosperm-reproduction:seed}{graines} étaient prises par l'envahisseur, et les géospizes à bec
moyen qui les lui disputaient furent perdants~: le sens de la sélection
était fixé par le compétiteur.
\end{solution}

\begin{solution}{exo:b2:speciation:12}
Faites par accident~: le modèle à deux locus montre l'isolement
apparaissant comme un sous-produit de substitutions fixées pour d'autres
raisons, ou par dérive, et la géographie --- une barrière, une fondation
--- fournit la séparation qui permet à l'accident de s'accumuler.
Maintenues par la sélection~: là où les formes se rencontrent, la
sélection contre les hybrides désavantagés construit des barrières
prézygotiques (renforcement), et la compétition déplace les caractères
pour que les deux puissent coexister. La formule omet la \omterm{prop:b2:speciation:modes}{spéciation
sympatrique} par sélection disruptive, où la sélection fait l'espèce
autant qu'elle la maintient, et la \omterm{def:b2:genome-diversification:polyploidy}{polyploïdie}, où l'accident est toute
l'histoire.
\end{solution}

\begin{solution}{pb:b2:speciation:1}
\textbf{1.} Espèce 1~: $N_1 + 0.9N_2 = 1200$, ordonnées à l'origine $N_1 = 1200$
et $N_2 = 1333$. Espèce 2~: $N_2 + 0.3N_1 = 400$, ordonnées à l'origine $N_2 =
400$ et $N_1 = 1333$.
\textbf{2.} $\beta K_1 = 360 < K_2 = 400$~: l'espèce 2 envahit~;
$\alpha K_2 = 360 < K_1 = 1200$~: l'espèce 1 envahit. Coexistence
stable.
\textbf{3.} $N_1^{*} = (1200 - 360)/(1 - 0.27) = 1151$~; $N_2^{*} =
(400 - 360)/0.73 = 55$.
\textbf{4.} $\mathrm{d}N_1/\mathrm{d}t = 0.5\times 1200\,(1 - 1218/1200)
= -9$ par an~: le résident, à sa capacité de charge, est entraîné dans
un déclin par les nouveaux venus.
\textbf{5.} $K_1 = 600$, $K_2 = 200$~: $N_1^{*} = (600 - 180)/0.73 =
575$, $N_2^{*} = (200 - 180)/0.73 = 27$~: les deux effectifs sont divisés
par deux, et l'envahisseur, avec 27 oiseaux, est proche de l'extinction
--- c'est pendant la sécheresse que la compétition se fait sentir.
\textbf{6.} L'envahisseur est l'oiseau le plus grand, au bec le plus
gros~: il prend les grosses \omterm{def:b2:angiosperm-reproduction:seed}{graines} que le résident utilise aussi, et
les prend plus vite, tandis que les petites \omterm{def:b2:angiosperm-reproduction:seed}{graines} du résident lui sont
de peu d'utilité.
\textbf{7.} $S = 9.0 - 9.5 = \qty{-0.5}{mm}$~; $R = 0.7\times(-0.5) =
\qty{-0.35}{mm}$.
\textbf{8.} $9.5 - 0.35 = \qty{9.15}{mm}$.
\textbf{9.} $N_1^{*} = (1200 - 200)/(1 - 0.15) = 1176$~; $N_2^{*} =
(400 - 360)/0.85 = 47$.
\textbf{10.} L'ordonnée à l'origine sur l'axe $N_2$ de l'isocline de
l'espèce 1 est passée de
$K_1/0.9 = 1333$ à $K_1/0.5 = 2400$~: la droite a pivoté en s'éloignant
de celle de l'envahisseur, le résident est désormais moins affecté par
chaque envahisseur, et l'équilibre s'est rapproché du $K$ du résident.
\textbf{11.} $2/0.35 \approx 6$ ans. Ce n'est pas le cas parce qu'une
sélection de ce type ne s'exerce que les années de sécheresse, s'inverse
les années humides où les petites \omterm{def:b2:angiosperm-reproduction:seed}{graines} abondent (ou lorsque
l'envahisseur est rare), et parce que la variance additive s'épuiserait.
\textbf{12.} Deux millions d'années représentent environ $10^{6}$
générations~; un déplacement de \qty{0.35}{mm} par génération soutenu
pendant ne serait-ce qu'un millième d'entre elles changerait le bec de
plusieurs centaines de millimètres. La sélection est forte mais
inconstante~; le changement directionnel soutenu est l'exception rare, et
le changement net sur les temps géologiques est un infime reliquat
d'épisodes importants et opposés.
\textbf{13.} Le chant, appris du père et utilisé par les femelles dans
le choix du partenaire, la forme du bec servant elle-même de signal. Les
hybrides aux becs intermédiaires ne traitent au mieux ni les grosses ni
les petites \omterm{def:b2:angiosperm-reproduction:seed}{graines} et meurent de faim les premiers lors des sécheresses.
\textbf{14.} $w = 0.5\sqrt{8/0.2} = \qty{3.2}{km}$.
\textbf{15.} L'île est plus étroite que ne le serait la zone~: aucun
cline spatial ne peut se former~; l'hybridation, là où elle se produit,
touche toute la population, et seuls les accouplements homogames
maintiennent les deux espèces séparées.
\textbf{16.} Les femelles qui ne s'accouplent qu'avec des mâles de leur
propre chant et de leur propre bec laissent davantage de petits-enfants
que celles qui s'hybrident~; la préférence et le signal (chant, bec)
deviennent plus distincts là où les deux espèces vivent ensemble que là
où elles vivent seules.
\textbf{17.} $0.005\times 15\times 15 = 1.1$ incompatibilité~: une en
moyenne, avec des hybrides probablement moins aptes mais non stériles.
Pas encore des espèces selon le critère biologique~; en voie de le
devenir.
\textbf{18.} Biologique~: pas encore (hybrides presque aussi aptes).
Morphologique~: peut-être, si les becs ou le plumage ont divergé
visiblement. Phylogénétique~: oui, si chaque population présente une
différence diagnostique fixée. Les concepts sont en désaccord précisément
pendant le processus de spéciation, ce qui fait leur utilité.
\textbf{19.} $1/(5\times 10^{6}) = 2\times 10^{-7}$ par espèce et par
an~; pour $10^{7}$ espèces, deux extinctions par an.
\textbf{20.} De 200 à 2000 espèces par an.
\textbf{21.} $0.75\times 10^{7}/10^{5} = 75$ espèces par an --- moins que
le taux actuel, même à sa valeur basse. L'épisode actuel, s'il se
prolonge, est une \omterm{prop:b2:speciation:tempo}{extinction de masse} plus rapide que les cinq grandes.
\textbf{22.} Les niches libérées par les disparus sont vides de
compétiteurs, si bien que tout variant survivant capable d'en exploiter
une est favorisé, et les survivants se diversifient pour les occuper.
\textbf{23.} \omterm{prop:b2:speciation:tempo}{Équilibre ponctué}~; \omterm{prop:b2:speciation:modes}{spéciation allopatrique} (péripatrique)
dans une petite population isolée dont le descendant envahit ensuite
l'aire de la population principale.
\textbf{24.} $1.04^{n} = 2$~: $n = \ln 2/\ln 1.04 = 18$ générations.
Les archives montrent des millions d'années parce que la sélection
s'inverse, que la \omterm{prop:b2:population-genetics:kinds}{sélection stabilisante} domine la plupart du temps, et
que le changement net est la petite différence entre des épisodes
importants et opposés.
\textbf{25.} Avant le déplacement~: coexistence stable en $N_1^{*} =
1151$, $N_2^{*} = 55$~; après~: 1176 et 47. Déplacement du bec de
\qty{-0.35}{mm} en une génération. Largeur de la \omterm{prop:b2:speciation:reinforcement}{zone hybride} \qty{3.2}{km}.
\end{solution}
